\documentclass[11pt,a4paper]{article}
\usepackage[etiket=en]{gtu-el-yazisi}
\begin{document}
\notbilgi{Analog Haberleşme}{Uygulama 1: Systems, Fourier Series \& Transform}{Analog-Uygulama-1.pdf}

\sayfa{1}
\begin{ornek}
  Classify these system into linear and non linear
  \begin{enumerate}[label=\alph*),nosep,leftmargin=*]
    \item $y(t)=2x(t)-3$
    \item $y(t)=e^{-t}x(t)$
  \end{enumerate}
\end{ornek}
Linearity Cond:
\[
  \left.\begin{aligned}
    x_1(t)&\xrightarrow{T[\cdot]}y_1(t)\\
    x_2(t)&\xrightarrow{T[\cdot]}y_2(t)
  \end{aligned}\right\}
  \quad T\bigl[ax_1(t)+bx_2(t)\bigr]\overset{?}{=}a\underbrace{T[x_1(t)]}_{y_1(t)}+b\underbrace{T[x_2(t)]}_{y_2(t)}
\]
a)
\begin{align*}
  x_1(t)&\to y_1(t)=2x_1(t)-3\\
  x_2(t)&\to y_2(t)=2x_2(t)-3\\
  x_3(t)&=ax_1(t)+bx_2(t)\\
  &\to y_3(t)=2x_3(t)-3=2ax_1(t)+2bx_2(t)-3\\
  &\overset{?}{\neq}a\bigl(2x_1(t)-3\bigr)+b\bigl(2x_2(t)-3\bigr)
\end{align*}
\vurgu[pembe]{non lineer!}

b) $y(t)=e^{-t}x(t)$
\begin{align*}
  x_1(t)&\to y_1(t)=e^{-t}x_1(t)\\
  x_2(t)&\to y_2(t)=e^{-t}x_2(t)\\
  x_3(t)&=ax_1(t)+bx_2(t)\\
  &\to y_3(t)=e^{-t}x_3(t)=e^{-t}ax_1(t)+e^{-t}bx_2(t)\\
  &\overset{?}{=}a\underbrace{y_1(t)}_{e^{-t}x_1(t)}+b\underbrace{y_2(t)}_{e^{-t}x_2(t)}\quad\checkmark
\end{align*}
\vurgu[pembe]{lineer system}

\begin{ornek}
  Determine whether the following systems are time variant or time invariant
  \begin{enumerate}[label=\alph*),nosep,leftmargin=*]
    \item $y(t)=2x(t)+3$
    \item $y(t)=(t+2)\,x(t)$
  \end{enumerate}
\end{ornek}
Time invariant cond:
\[ x(t)\to y(t),\qquad x(t-t_0)\to y(t-t_0) \]
a)
\[
  \left.\begin{aligned}
    x_1(t)&\to\overset{2x_1(t)+3}{y_1(t)}\\
    x_2(t)&=x(t-t_0)
  \end{aligned}\right\}
  \quad y_2(t)\overset{?}{=}y_1(t-t_0)
\]
\[ 2x_2(t)+3=2x(t-t_0)+3\overset{?}{=}2x_1(t-t_0)+3\quad\checkmark \]
zamanda değişmez

\sayfa{2}
b)
\begin{align*}
  x_1(t)&\to y_1(t)=(t+2)\,x_1(t)\\
  x_2(t)=x_1(t-t_0)&\to y_2(t)=(t+2)\,x_2(t)\\
  y_2(t)&=(t+2)\,x_1(t-t_0)\\
  y_2(t)&\overset{?}{=}y_1(t-t_0)\\
  (t+2)\,x_1(t-t_0)&\overset{?}{\neq}(t-t_0+2)\,x_1(t-t_0)
\end{align*}
time variant
\begin{ornek}
  A system is causal if its input at any time $t_0$ depends on the input at times prior to $t_0$
\end{ornek}
\[ y(t_0)=T\bigl[x(t):t\le t_0\bigr] \]
A nessesary an sufficent condition for an LTI system to be causal is that \underline{its impulse response $h(t)$,} \textcolor{kirmizikalem}{[the output when the input is $\delta(t)$]} must be a causal siggnal for $t<0$, we must have $h(t)=0$

$h(t)=\Pi\!\left(\dfrac{t-3}{6}\right)$, LTI system is causal?
\begin{center}
  \begin{tikzpicture}[ax/.style={-{Stealth[length=2mm]},line width=0.6pt}]
    % (1) A Pi(t/tau)
    \node[kirmizikalem] at (0,2.0) {$A\,\Pi(t/\tau)$};
    \draw[ax] (-1.7,0) -- (1.9,0);
    \draw[ax] (0,-0.1) -- (0,1.6);
    \node[right] at (0,1.45) {$A$};
    \draw[kirmizikalem,line width=0.9pt] (-1,0) -- (-1,1.1) -- (1,1.1) -- (1,0);
    \node[kirmizikalem,below,font=\small] at (-1,0) {$-\tau/2$};
    \node[kirmizikalem,below,font=\small] at (1,0) {$\tau/2$};
    % (2) Pi(t/6)
    \begin{scope}[shift={(5.6,0)}]
      \node at (0,2.0) {$\Pi(t/6)$};
      \draw[ax] (-1.7,0) -- (1.9,0);
      \draw[ax] (0,-0.1) -- (0,1.6);
      \draw[line width=0.9pt] (-1.05,0) -- (-1.05,1.1) -- (1.05,1.1) -- (1.05,0);
      \node[below,font=\small] at (-1.05,0) {$-3$};
      \node[below,font=\small] at (1.05,0) {$3$};
    \end{scope}
    % (3) h(t)
    \begin{scope}[shift={(10.6,0)}]
      \node at (0.8,2.0) {$h(t)=\Pi\!\left(\frac{t-3}{6}\right)$};
      \draw[ax] (-0.5,0) -- (3.0,0) node[right]{$t$};
      \draw[ax] (0,-0.1) -- (0,1.6);
      \draw[line width=0.9pt] (0,0) -- (0,1.1) -- (2.1,1.1) -- (2.1,0);
      \node[left,font=\small] at (0,1.1) {$1$};
      \node[below,font=\small] at (0,0) {$0$};
      \node[below,font=\small] at (2.1,0) {$6$};
      \draw[->] (0.6,-0.1) -- (0.6,-0.6);
      \node[below,font=\small] at (0.6,-0.6) {causal};
    \end{scope}
  \end{tikzpicture}
\end{center}
\begin{ornek}
  Determine the Fourier Series Expansion of the following signals.
\end{ornek}

\sayfa{3}
\[ \text{a)}\quad x(t)=\sum_{n=-\infty}^{\infty}\Lambda(t-n)\,u(t-n) \]
\[ \text{b)}\quad x(t)=\sum_{n=-\infty}^{\infty}(-1)^n\,\delta(t-nT) \]
\section{Fourier Seri Açılımları}
Trigonometrik FS:
\[ x(t)=a_0+\sum_{n=1}^{\infty}\bigl(a_n\cos(n\omega_0t)+b_n\sin(n\omega_0t)\bigr) \]
\[
  \begin{aligned}
    a_0&=\frac{1}{T_0}\int_{T_0}x(t)\,dt,\qquad a_n=\frac{2}{T_0}\int_{T_0}x(t)\cos(n\omega_0t)\,dt\\
    b_n&=\frac{2}{T_0}\int_{T_0}x(t)\sin(n\omega_0t)\,dt
  \end{aligned}
\]
kompleks ExpFS:
\[ x(t)=\sum_{n=-\infty}^{\infty}x_n\,e^{jn\omega_0t},\qquad x_n=\frac{1}{T_0}\int_{T_0}x(t)\,e^{-jn\omega_0t}\,dt \]
\formulkutu[ikisi arasındaki \okunamadi{bağıntı}]{x_n=\frac{a_n-jb_n}{2}}
a) önce $n=0$ ve $n=1$ için çizelim
\[ n=0\ \to\ \Lambda(t)\,u(t) \]
\begin{center}
  \begin{tikzpicture}[ax/.style={-{Stealth[length=2mm]},line width=0.6pt},cv/.style={line width=0.9pt}]
    \begin{scope}
      \draw[ax] (-1.5,0) -- (1.7,0) node[right]{$t$};
      \draw[ax] (0,-0.05) -- (0,1.5);
      \draw[cv] (-1,0) -- (0,1) -- (1,0);
      \node[below,font=\small] at (-1,0) {$-1$};
      \node[below,font=\small] at (1,0) {$1$};
    \end{scope}
    \begin{scope}[shift={(5.3,0)}]
      \draw[ax] (-1.5,0) -- (1.9,0) node[right]{$t$};
      \draw[ax] (0,-0.05) -- (0,1.5);
      \draw[cv] (0,1) -- (1.7,1);
      \node[left,font=\small] at (0,1) {$1$};
    \end{scope}
    \begin{scope}[shift={(10.6,0)}]
      \draw[ax] (-1.5,0) -- (1.7,0);
      \draw[ax] (0,-0.05) -- (0,1.5);
      \draw[cv] (0,0) -- (0,1) -- (1,0);
      \node[left,font=\small] at (0,1) {$1$};
      \node[below,font=\small] at (1,0) {$1$};
    \end{scope}
  \end{tikzpicture}
\end{center}
\[ n=1\ \to\ \Lambda(t-1)\,u(t-1) \]
\begin{center}
  \begin{tikzpicture}[ax/.style={-{Stealth[length=2mm]},line width=0.6pt},cv/.style={line width=0.9pt}]
    \begin{scope}[kirmizikalem]
      \node[anchor=east] at (-1.6,0.9) {$A\,\Lambda(t/2)$};
      \draw[ax] (-2.3,0) -- (2.4,0) node[right]{$t$};
      \draw[ax] (0,-0.05) -- (0,1.6);
      \draw[cv] (-2,0) -- (0,1.2) -- (2,0);
      \node[right] at (0,1.35) {$A$};
      \node[below,font=\small] at (-2,0) {$-2$};
      \node[below,font=\small] at (2,0) {$2$};
    \end{scope}
    \begin{scope}[shift={(8,0)}]
      \draw[ax] (-1.2,0) -- (3.4,0) node[right]{$t$};
      \draw[ax] (0,-0.05) -- (0,1.5);
      \draw[dashed,line width=0.4pt] (0,1) -- (3.0,1);
      \draw[cv] (1,0) -- (1,1) -- (2,0);
      \node[left,font=\small] at (0,1) {$1$};
      \node[below,font=\small] at (1,0) {$1$};
      \node[below,font=\small] at (2,0) {$2$};
    \end{scope}
  \end{tikzpicture}
\end{center}
\[ \sum_{n=-\infty}^{\infty}\Lambda(t-n)\,u(t-n) \]
\begin{center}
  \begin{tikzpicture}[ax/.style={-{Stealth[length=2mm]},line width=0.6pt},cv/.style={line width=0.9pt}]
    \begin{scope}[x=0.85cm,y=1.6cm]
      \draw[ax] (-4.6,0) -- (4.9,0) node[right]{$t$};
      \draw[ax] (0,-0.03) -- (0,1.3) node[right]{$x(t)$};
      \draw[dashed,line width=0.4pt] (-4.3,1) -- (4.6,1);
      \foreach \n in {-4,-3,...,3} \draw[cv] (\n,1) -- (\n+1,0);
      \foreach \n in {-4,-3,...,4} \draw[dashed,line width=0.4pt] (\n,0) -- (\n,1);
      \foreach \n in {-4,-3,...,4} \node[below,font=\small] at (\n,0) {$\n$};
    \end{scope}
  \end{tikzpicture}
\end{center}

\sayfa{4}
kompleks FS
\[
  \begin{aligned}
    x_0&=\frac{1}{T_0}\int_{T_0}x(t)\,dt\\
    &=\frac{1}{T_0}\int x(t)\,dt\\
    &=\int_0^1(1-t)\,dt=\frac{1}{2}
  \end{aligned}
\]
\underline{$T_0=1$}
\[ y-y_0=\underset{-1}{\overset{\text{doğrunun eğimi}}{m}}\,(x-x_0) \]
\[ y-1=(-1)\,x \]
\[ y(t)=1-t \]
\underline{$T_0=1$}
\[
  n\neq0\ \to\ x_n=\int_0^1x(t)\,e^{-jn2\pi t}=\int_0^1(1-t)\,e^{-jn2\pi t}\,dt
  =\int_0^1e^{-jn2\pi t}\,dt\ \underbrace{-\int_0^1t\,e^{-jn2\pi t}\,dt}_{I_1}
\]
\[ \omega_0=\frac{2\pi}{T_0} \]
\[ \downarrow\ \ -\frac{1}{jn2\pi}\,e^{-jn2\pi t}\Big|_0^1=\frac{-1}{jn2\pi}\Bigl[e^{-jn2\pi}-e^{0}\Bigr] \]
\[ I_1=\int_0^1t\,e^{-jn2\pi t}\,dt \]
\formulkutu[kısmi integrasyon]{\left.uv\right|_a^b-\int_a^b v\,du}
\[
  \begin{aligned}
    u&=t & e^{-jn2\pi t}\,dt&=dv\\
    du&=dt & \frac{-1}{jn2\pi}\,e^{-jn2\pi t}&=v
  \end{aligned}
\]
\[ =\left(t\cdot\frac{-1}{jn2\pi}\,e^{-jn2\pi t}\right)\Bigg|_0^1-\int_0^1-\frac{1}{jn2\pi}\,e^{-jn2\pi t}\,dt \]
\[ =-\frac{1}{jn2\pi}\,e^{-jn2\pi}+\frac{1}{n^24\pi^2}\,e^{-jn2\pi}-\frac{1}{n^24\pi^2} \]
\[ x_n=\frac{1}{jn2\pi}+\frac{1}{n^24\pi^2}-\frac{1}{n^24\pi^2}\,e^{-jn2\pi} \]
\[ =\frac{1}{jn2\pi}+\frac{1}{n^24\pi^2}\Bigl[1-\underbrace{e^{-jn2\pi}}_{1}\Bigr]=\frac{1}{jn2\pi} \]
\[ \text{Üstel FS:}\quad\kutula{x(t)=\frac{1}{2}+\sum_{n\neq0}\frac{1}{jn2\pi}\,e^{jn2\pi t}} \]

\sayfa{5}
\[ x_n=\frac{a_n-jb_n}{2}=\overset{\text{$j$ ile çarpalım}}{\frac{1}{jn2\pi}}=-\frac{j}{n2\pi} \]
\[ a_n=0,\qquad b_n=\frac{1}{n\pi} \]
\[ x(t)=\frac{1}{2}+\sum_{n=1}^{\infty}\frac{1}{\pi n}\sin(2\pi nt) \]
b)
\[ \begin{aligned} n=0&\to\delta(t)\\ n=1&\to-\delta(t-T) \end{aligned} \]
\begin{center}
  \begin{tikzpicture}[ax/.style={-{Stealth[length=2mm]},line width=0.6pt},ar/.style={-{Stealth[length=2mm]},line width=1pt}]
    \draw[ax] (-1.5,0) -- (2.0,0) node[right]{$t$};
    \draw[ar] (0,0) -- (0,1.0);
    \node[below,font=\small] at (0,0) {$0$};
    \begin{scope}[shift={(6.5,0)}]
      \draw[ax] (-1.2,0) -- (2.4,0) node[right]{$t$};
      \draw[ar] (1.2,0) -- (1.2,-1.0);
      \node[above,font=\small] at (1.2,0) {$T$};
      \node[below,font=\small] at (0,0) {$0$};
    \end{scope}
  \end{tikzpicture}
\end{center}
c)
\begin{center}
  \begin{tikzpicture}[ax/.style={-{Stealth[length=2mm]},line width=0.6pt},ar/.style={-{Stealth[length=2mm]},line width=1pt}]
    \draw[ax] (-3.7,0) -- (4.9,0) node[right]{$t$};
    \draw[line width=0.6pt] (0,-1.3) -- (0,1.5);
    \foreach \n/\lab in {-3/-3T,-2/-2T,-1/-T,0/0,1/T,2/2T,3/3T,4/4T} {
      \ifodd\n
        \draw[ar] (\n,0) -- (\n,-0.9); \node[right,font=\small] at (\n,-0.4) {$\lab$};
      \else
        \draw[ar] (\n,0) -- (\n,0.9); \node[below,font=\small] at (\n,0) {$\lab$};
      \fi}
    \node[right,font=\small] at (0,0.9) {$1$};
  \end{tikzpicture}
\end{center}
\textcolor{kirmizikalem}{Dürtü katarının FS açılımı}
\[ x(t)=\sum_n\delta(t-nT_0)=\frac{1}{T_0}\sum_n e^{jn\omega_0t} \]
\[ x(t)=\underbrace{\sum_n\delta(t-n2T)}_{}\;-\;\underbrace{\sum_n\delta(t-T-n2T)}_{} \]
\begin{center}
  \begin{tikzpicture}[ax/.style={-{Stealth[length=2mm]},line width=0.6pt},ar/.style={-{Stealth[length=2mm]},line width=1pt}]
    \begin{scope}[x=0.7cm]
      \draw[ax] (-4.8,0) -- (5.2,0);
      \foreach \n/\lab in {-4/-4T,-2/-2T,0/0,2/2T,4/4T} {
        \draw[ar] (\n,0) -- (\n,0.9); \node[below,font=\small] at (\n,0) {$\lab$};}
    \end{scope}
    \begin{scope}[shift={(8.2,0)},x=0.7cm]
      \draw[ax] (-4.8,0) -- (5.2,0);
      \foreach \n/\lab in {-3/-3T,-1/-T,1/T,3/3T} {
        \draw[ar] (\n,0) -- (\n,0.9); \node[below,font=\small] at (\n,0) {$\lab$};}
    \end{scope}
  \end{tikzpicture}
\end{center}
\[ x(t)=\frac{1}{2T}\sum_n e^{jn\frac{2\pi}{2T}t}-\frac{1}{2T}\sum_n e^{jn\frac{2\pi}{2T}(t-T)} \]
\[ =\frac{1}{2T}\sum_n\Bigl\{e^{jn\frac{\pi}{T}t}-e^{jn\frac{\pi}{T}t}\,\underbrace{e^{-jn\frac{\pi}{T}\cdot T}}_{(-1)^n}\Bigr\} \]
\[ \textcolor{kirmizikalem}{e^{-jn\pi}=\cos(-n\pi)+j\sin(-n\pi)=(-1)^n} \]
\[ =\frac{1}{2T}\sum_n e^{jn\frac{\pi}{T}\cdot t}\,\bigl(1-(-1)^n\bigr) \]
\[ x_n=\frac{1}{2T}\bigl(1-(-1)^n\bigr)=\begin{cases}\dfrac{1}{T}&n\ \text{tek}\\[2mm] 0&n\ \text{çift}\end{cases} \]

\sayfa{6}
\formulkutu{\Rightarrow\ x(t)=\sum_n x_n\,e^{jn\frac{\pi}{T}\cdot t},\qquad x_n=\begin{cases}1/T,&n\ \text{tek}\\ 0,&n\ \text{çift}\end{cases}}
\begin{ornek}
  Determine the FS expansion of the following periodic signal
\end{ornek}
\elcizimi[0.65]{sekiller/uygulama-notlari-Analog-Uygulama-1-s6-1.png}[6]
\[ y=\frac{1}{T}\,x,\qquad x(t)=\frac{t}{T} \]
\[ x_n=\frac{1}{T_0}\int_{T_0}x(t)\,e^{-jn\omega_0t}\,dt \]
\[ T_0=2T\ \to\ \omega_0=\frac{2\pi}{2T}=\frac{\pi}{T} \]
\noindent\begin{minipage}{\linewidth}
\[
  \begin{aligned}
    x_n&=\frac{1}{2T}\int_{-T}^{T}\frac{t}{T}\,e^{-jn\frac{\pi}{T}}\,dt\\
    &=\frac{1}{2T^2}\int_{-T}^{T}e^{-jn\frac{\pi t}{T}}\,dt\text{\footnotemark}
  \end{aligned}
\]
\footnotetext{Claude notu: birinci satırda üstelde, ikinci satırda ise integrandda $t$ çarpanı yazılmamış; sonraki adımlara göre ($u=t$) muhtemelen $x_n=\frac{1}{2T^2}\int_{-T}^{T}t\,e^{-jn\frac{\pi t}{T}}\,dt$ olmalı.}
\end{minipage}
\[
  \begin{aligned}
    u&=t & e^{-jn\frac{\pi t}{T}}\,dt&=dv\\
    du&=dt & -\frac{1}{jn\frac{\pi}{T}}\,e^{-jn\frac{\pi t}{T}}&=v
  \end{aligned}
\]
\[
  \begin{aligned}
    &=\frac{1}{2T^2}\Biggl[\Bigl(t\cdot\frac{-1}{jn\frac{\pi}{T}}\,e^{-jn\frac{\pi t}{T}}\Bigr)\Bigg|_{-T}^{T}\\
    &\qquad-\int_{-T}^{T}-\frac{1}{jn\frac{\pi}{T}}\,e^{-jn\frac{\pi t}{T}}\,dt\Biggr]
  \end{aligned}
\]
Sonuca giderken $e^{jn\pi}=(-1)^n$

$x_0$'ı ayrı olarak bul ödev

\sayfa{7}
\begin{ornek}
  Classify the following signal into periodic or non periodic
  \[ x(t)=2\sum_{n}\Lambda\!\left(\frac{t-4n}{2}\right)-\sum_{n}\Lambda(t-4n) \]
\end{ornek}
\[ n=0\ \to\ x_0(t)=x_1(t)=2\Lambda\!\left(\frac{t}{2}\right)-\Lambda(t) \]
\elcizimi[0.65]{sekiller/uygulama-notlari-Analog-Uygulama-1-s7-1.png}[7]

bunların farkı
\elcizimi[0.4]{sekiller/uygulama-notlari-Analog-Uygulama-1-s7-2.png}[7]
\elcizimi[0.85]{sekiller/uygulama-notlari-Analog-Uygulama-1-s7-3.png}[7]
\textcolor{kirmizikalem}{$T_0=4$ ile periyodiktir}

\begin{ornek}
  Find the Fourier Transform of the following signal
\end{ornek}
\elcizimi[0.65]{sekiller/uygulama-notlari-Analog-Uygulama-1-s7-4.png}[7]
\[ X(f)=\,? \]
\elcizimi[0.9]{sekiller/uygulama-notlari-Analog-Uygulama-1-s7-5.png}[7]

\sayfa{8}
\[ x(t)=\sin(2\pi f_0t)\cdot\Pi\!\left(\frac{t-\frac{1}{4f_0}}{\frac{1}{2f_0}}\right)-\sin(2\pi f_0t)\,\Pi\!\left(\frac{t+\frac{1}{4f_0}}{\frac{1}{2f_0}}\right) \]
\[ x(t)=\underbrace{\sin(2\pi f_0t)}_{x_1(t)}\Bigl[\underbrace{\Pi\Bigl(2f_0\bigl(t-\tfrac{1}{4f_0}\bigr)\Bigr)-\Pi\Bigl(2f_0\bigl(t+\tfrac{1}{4f_0}\bigr)\Bigr)}_{x_2(t)}\Bigr] \]
\[
  \begin{aligned}
    x(t)&=x_1(t)\cdot x_2(t)\\
    \downarrow\\
    X(f)&=X_1(f)\circledast X_2(f)
  \end{aligned}
\]
\[ \textcolor{kirmizikalem}{\sin(2\pi f_0t)\xrightarrow{\mathcal{F}}-\frac{1}{2j}\Bigl[\delta(f+f_0)-\delta(f-f_0)\Bigr]\quad\leftarrow X_1(f)} \]
\[ \textcolor{kirmizikalem}{A\,\Pi\!\left(\frac{t}{\tau}\right)\xrightarrow{\mathcal{F}}A\,\tau\,\mathrm{sinc}(f\tau)} \]
\[ \textcolor{kirmizikalem}{x(t-t_0)\to X(f)\,e^{-j2\pi ft_0}} \]
\[
  x_2(t)=\Pi\Bigl(\underbrace{2f_0}_{\tau=\frac{1}{2f_0}}\bigl(t-\underbrace{\tfrac{1}{4f_0}}_{t_0=\frac{1}{4f_0}}\bigr)\Bigr)-\Pi\Bigl(2f_0\bigl(t+\underbrace{\tfrac{1}{4f_0}}_{t_0=-\frac{1}{4f_0}}\bigr)\Bigr)
\]
\[
  \begin{aligned}
    X_2(f)&=\frac{1}{2f_0}\,\mathrm{sinc}\Bigl(f\cdot\frac{1}{2f_0}\Bigr)e^{-j2\pi f\cdot\frac{1}{4f_0}}\\
    &\quad-\frac{1}{2f_0}\,\mathrm{sinc}\Bigl(f\cdot\frac{1}{2f_0}\Bigr)\cdot e^{j2\pi f\cdot\frac{1}{4f_0}}\\
    &=\frac{j}{f_0}\cdot\mathrm{sinc}\Bigl(\frac{f}{2f_0}\Bigr)\sin\Bigl(\frac{\pi f}{2f_0}\Bigr)
  \end{aligned}
\]
\[ \textcolor{kirmizikalem}{X(f)=X_2(f)\circledast\Bigl[\frac{1}{2j}\bigl(\delta(f-f_0)-\delta(f+f_0)\bigr)\Bigr]} \]
\[ \textcolor{kirmizikalem}{=\frac{1}{2j}\underbrace{X_2(f)\circledast\delta(f-f_0)}_{X_2(f-f_0)}-\frac{1}{2j}\underbrace{X_2(f)\circledast\delta(f+f_0)}_{X_2(f+f_0)}} \]
\[ \textcolor{kirmizikalem}{\begin{aligned} x(t)\,\delta(t-t_0)&=x(t_0)\cdot\delta(t-t_0)\\ x(t)\circledast\delta(t-t_0)&=x(t-t_0) \end{aligned}} \]

\sayfa{9}
\begin{ornek}
  \[ x(t)=4\Lambda\!\left(\frac{t}{4}\right)\cos(2\pi f_0t) \]
  Find the Fourier Transform
\end{ornek}
\textcolor{kirmizikalem}{Reminder:}
\[
  \begin{aligned}
    A\Lambda(t/\tau)&\to A\tau\,\mathrm{sinc}^2(f\tau)\\
    \cos(2\pi f_0t)&\to\frac{1}{2}\delta(f-f_0)+\frac{1}{2}\delta(f+f_0)\\
    x(t)\cos(2\pi f_0t)&\to\frac{1}{2}\bigl[X(f-f_0)+X(f+f_0)\bigr]
  \end{aligned}
\]

\begin{ornek}
  \[ x(t)=\Lambda(2t+3)-\Lambda(3t-2) \]
  Find the fourier transform
\end{ornek}
Reminder
\[
  \begin{aligned}
    A\Lambda(t/\tau)&\to A\tau\,\mathrm{sinc}^2(f\tau)\\
    x(t-t_0)&\to X(f)\,e^{-j2\pi ft_0}
  \end{aligned}
\]
\[ x(t)=\underbrace{\Lambda\!\left(\frac{t+3/2}{1/2}\right)}_{\substack{A=1\\ \tau=1/2\\ t_0=-3/2}}-\underbrace{\Lambda\!\left(\frac{t-2/3}{1/3}\right)}_{\substack{A=1\\ \tau=1/3\\ t_0=2/3}} \]
\[ X(f)=\frac{1}{2}\,\mathrm{sinc}^2\!\left(\frac{f}{2}\right)e^{j2\pi f\frac{3}{2}}-\frac{1}{3}\,\mathrm{sinc}^2\!\left(\frac{f}{3}\right)e^{-j2\pi f\frac{2}{3}} \]

\begin{ornek}
  Fourier transform of a signal is
\end{ornek}
\elcizimi[0.3]{sekiller/uygulama-notlari-Analog-Uygulama-1-s9-1.png}[9]
Determine and sketch the fourier transform of the signal
\[
  x_1(t)=\mathord{\ooalign{$-x(t)$\cr\hidewidth\textcolor{kirmizikalem}{\rule[0.6ex]{2.4em}{0.6pt}}\hidewidth\cr}}
  +x(t)\underbrace{\cos(2000\pi t)}_{\displaystyle\textcolor{kirmizikalem}{2\pi f_0t}}
  +\overset{\textcolor{kirmizikalem}{\uparrow\,x(t)}}{2x(t)}\underbrace{\cos^2(3000\pi t)}_{\displaystyle\textcolor{kirmizikalem}{\frac{1}{2}\bigl[1+\cos(6000\pi t)\bigr]}}
\]
\textit{(Not: $-x(t)$ terimi kırmızı kalemle çizilerek iptal edilmiş.)}

\sayfa{10}
\[
  \begin{aligned}
    X_1(f)&=\frac{1}{2}X(f-1000)+\frac{1}{2}X(f+1000)+\frac{1}{2}X(f-3000)+{}\\
    &\quad\frac{1}{2}X(f+3000)
  \end{aligned}
\]
\[ \cos^2\alpha=\frac{1}{2}\bigl(1+\cos2\alpha\bigr) \]

\begin{ornek}
  \[ x(t)=t\,\mathrm{sinc}(t) \]
  Find the fourier transform
\end{ornek}
\textcolor{kirmizikalem}{1. yol}
\[ \mathrm{sinc}(t)=\frac{\sin(\pi t)}{\pi t}\qquad\textcolor{kirmizikalem}{\longrightarrow\ 2\pi f_0t} \]
\[ x(t)=\frac{1}{\pi}\sin(\pi t)\qquad\begin{aligned}[t]2\pi f_0t&=\pi t\\ f_0&=\,?\,=\frac{1}{2}\end{aligned} \]
\[ X(f)=\frac{1}{2j\pi}\,\delta\!\left(f-\frac{1}{2}\right)-\frac{1}{2j\pi}\,\delta\!\left(f+\frac{1}{2}\right) \]

\textcolor{kirmizikalem}{2. yol \quad Reminder:}
\[ \textcolor{kirmizikalem}{t\,x(t)\to\frac{j}{2\pi}\,\frac{dX(f)}{df}} \]
\[ x_1(t)=\mathrm{sinc}(t)\to X_1(f)=\Pi(f) \]
\elcizimi[0.3]{sekiller/uygulama-notlari-Analog-Uygulama-1-s10-1.png}[10]
\[ X_1(f)=U\!\left(f+\frac{1}{2}\right)-U\!\left(f-\frac{1}{2}\right) \]
\[ \frac{dX_1(f)}{df}=\delta\!\left(f+\frac{1}{2}\right)-\delta\!\left(f-\frac{1}{2}\right) \]
\[ X(f)=\frac{j}{2\pi}\Bigl[\delta\!\left(f+\frac{1}{2}\right)-\delta\!\left(f-\frac{1}{2}\right)\Bigr] \]

\begin{ornek}
  \[ x(t)=\mathrm{sinc}(t)\circledast\mathrm{sinc}(t)=\mathrm{sinc}(t) \]
  Prove.
\end{ornek}
\[ \textcolor{kirmizikalem}{\mathcal{F}\{x_1(t)\ast x_2(t)\}=X_1(f)\cdot X_2(f)} \]
\[
  \begin{array}{rcl}
    x_1(t)=\mathrm{sinc}(t)&\longrightarrow&X_1(f)=\Pi(f)\\
    x_2(t)=\mathrm{sinc}(t)&\longrightarrow&X_2(f)=\Pi(f)\\
    &&X(f)=\Pi(f)\\
    &&\downarrow\mathcal{F}^{-1}\\
    &&x(t)=\mathrm{sinc}(t)
  \end{array}
\]
\elcizimi[0.3]{sekiller/uygulama-notlari-Analog-Uygulama-1-s10-2.png}[10]

\sayfa{11}
\begin{ornek}
  \[ x(t)=\frac{1}{1+t^2}\qquad\qquad X(f)=\,? \]
\end{ornek}
\textcolor{kirmizikalem}{Reminder:}
\[ \textcolor{kirmizikalem}{e^{-\alpha|t|}\to\frac{2\alpha}{\alpha^2+(2\pi f)^2}}\ \to\ \frac{2\alpha}{4\pi^2}\Bigl[\frac{1}{\frac{\alpha^2}{4\pi^2}+f^2}\Bigr] \]
Dualite
\[
  \begin{aligned}
    x(t)&\to X(f)\\
    X(t)&\to x(-f)
  \end{aligned}
\]
\[ \frac{\alpha^2}{4\pi^2}=1\ \to\ \alpha^2=4\pi^2\ \Rightarrow\ \alpha=2\pi \]
\[ X(t)=\frac{2\cdot2\pi}{4\pi^2}\Bigl[\frac{1}{1+t^2}\Bigr]=\frac{1}{\pi}\,\frac{1}{1+t^2} \]
\[ x(-f)=e^{-\alpha|-f|}=e^{-\alpha|f|}=\pi\cdot e^{-\alpha|f|}=\pi\cdot e^{-2\pi|f|} \]

\begin{ornek}
  \[ X(f)=\,?\,=\int_{-\infty}^{\infty}x(t)\,e^{-j2\pi ft}\,dt \]
\end{ornek}
\elcizimi[0.6]{sekiller/uygulama-notlari-Analog-Uygulama-1-s11-1.png}[11]
\[ X(f)=\int_{-1}^{0}(t+1)\,e^{-j2\pi ft}\,dt+\int_{0}^{1}(t-1)\,e^{-j2\pi ft}\,dt \]
\[ \text{sonuç}\quad=\frac{j}{\pi f}\bigl[1-\mathrm{sinc}(2f)\bigr] \]
\[
  \begin{aligned}
    &=\underbrace{\int_{-1}^{0}t\,e^{-j2\pi ft}\,dt}_{\textcolor{kirmizikalem}{I_1}}+\underbrace{\int_{-1}^{0}e^{-j2\pi ft}\,dt}_{\displaystyle\frac{-1}{j2\pi f}\,e^{-j2\pi ft}\Big|_{-1}^{0}}\\
    &\quad+\underbrace{\int_{0}^{1}t\,e^{-j2\pi ft}\,dt}_{\textcolor{kirmizikalem}{I_2}}-\underbrace{\int_{0}^{1}e^{-j2\pi ft}\,dt}_{\displaystyle\frac{-1}{j2\pi f}\,e^{-j2\pi ft}\Big|_{0}^{1}}
  \end{aligned}
\]

\sayfa{12}
\textcolor{kirmizikalem}{kısmi integrasyon}
\[ \textcolor{kirmizikalem}{\int_a^b u\,dv=u\,v\Big|_a^b-\int_a^b v\,du} \]
\[ I_1=\int_{-1}^{0}t\,e^{-j2\pi ft}\,dt= \]
\[
  \begin{array}{ll}
    u=t & e^{-j2\pi ft}\,dt=dv\\
    du=dt & {}=\dfrac{-1}{j2\pi f}\,e^{-j2\pi ft}=v
  \end{array}
\]
\[ {}=\Bigl(t\,\frac{-1}{j2\pi f} \]
\end{document}
